\documentclass[11pt,a4paper]{article}

% --- packages ---
\usepackage[utf8]{inputenc}
\usepackage[T1]{fontenc}
\usepackage[spanish]{babel}
\usepackage{lmodern}
\usepackage{microtype}
\usepackage{booktabs}
\usepackage{array}
\usepackage{tabularx}
\usepackage{longtable}
\usepackage{enumitem}
\usepackage{geometry}
\usepackage{hyperref}
\usepackage{amsmath}
\usepackage{amssymb}
\usepackage{graphicx}
\usepackage{xcolor}

\geometry{margin=2.5cm}
\hypersetup{
  colorlinks=true,
  linkcolor=black,
  urlcolor=blue!60!black,
  citecolor=black,
  pdftitle={Metricas Expected en el Futbol}
}

% --- metadata ---
\title{
  Métricas \emph{Expected} en el Fútbol \\
  \large Marco conceptual, fuentes de datos y construcción
}
\author{Proyecto \textsc{Futbolstats}}
\date{\today}

% --- document ---
\begin{document}

\maketitle
\begin{abstract}
\noindent
Este documento describe el marco teórico de las métricas \emph{expected} (xG, xA, xT, xCards, xCorners, etc.), las fuentes de datos públicas que permiten construirlas y la implementación concreta de este pipeline en el proyecto \textsc{Futbolstats}. La motivación de fondo es entrenar modelos predictivos a nivel \emph{jugador-partido} sobre estadísticas individuales, no sobre resultados, y este texto formaliza qué se necesita para hacerlo bien.
\end{abstract}

\tableofcontents
\clearpage

\section{Introducción}
\label{sec:intro}

\subsection{Objetivo}

Este documento sirve como referencia técnica del marco conceptual detrás del proyecto \textsc{Futbolstats}. Su lectura no es necesaria para correr los scrapers, pero sí lo es para entender \emph{por qué} se eligen ciertas métricas como objetivo de modelado y \emph{qué datos} habilitan cada una.

\subsection{Contexto del proyecto}

El proyecto \textsc{Futbolstats} construye una base de datos estructurada de estadísticas de fútbol desde fuentes públicas (FBref, StatsBomb Open Data, Understat, ClubElo) con el objetivo de entrenar modelos predictivos de \emph{estadísticas individuales por partido}. Concretamente, se predicen variables como:

\begin{itemize}[noitemsep]
  \item tiros, tiros al arco, xG generado;
  \item pases intentados y completados, pases clave;
  \item entradas, intercepciones, despejes;
  \item faltas cometidas y recibidas, tarjetas amarillas y rojas;
  \item dribles intentados y completados.
\end{itemize}

\noindent
El enfoque es deliberadamente distinto al de la mayoría de literatura aplicada en fútbol, que se concentra en predecir \emph{resultados} (victoria/empate/derrota) o \emph{goles totales} por partido. Predecir resultados es un problema de baja dimensionalidad y alto ruido; predecir estadísticas individuales abre un espacio mucho más rico de modelado y permite evaluar componentes específicos del juego.

\subsection{Audiencia}

Este texto asume:

\begin{itemize}[noitemsep]
  \item familiaridad con regresión y validación cruzada temporal;
  \item conocimiento básico de fútbol (posiciones, fases de juego);
  \item interés en construir features desde datos brutos, no sólo consumir métricas pre-calculadas.
\end{itemize}

\subsection{Estructura del documento}

La Sección~\ref{sec:xg} introduce formalmente el concepto de \emph{Expected Goals} (xG) y explica por qué la idea se generaliza naturalmente a cualquier evento contable del juego. La Sección~\ref{sec:catalogo} presenta el catálogo de métricas \emph{expected} más utilizadas en la literatura, con definición, interpretación y observables necesarios para cada una. La Sección~\ref{sec:fuentes} cruza ese catálogo con las fuentes de datos públicas disponibles, indicando qué métrica se puede construir desde dónde y con qué cobertura. La Sección~\ref{sec:impl} describe la implementación concreta en este repositorio: scripts, schemas y flujo de procesamiento. La Sección~\ref{sec:limitaciones} discute las limitaciones del enfoque actual y posibles líneas de trabajo futuro.

\subsection{Convenciones de notación}

\begin{itemize}[noitemsep]
  \item $x_i$ denota un evento individual (un tiro, un pase, una falta).
  \item $\mathbf{c}_i$ denota el vector de contexto observado en el momento de $x_i$ (posición en cancha, parte del cuerpo, presión defensiva, etc.).
  \item $Y_i \in \{0,1\}$ es la realización binaria del evento (gol/no gol, tarjeta/no tarjeta, etc.).
  \item $\widehat{Y}_i = P(Y_i = 1 \mid \mathbf{c}_i)$ es la predicción del modelo, interpretada como probabilidad.
  \item $\mathrm{x}\textsc{M}$ con \textsc{M} en versalitas refiere a la métrica \emph{expected} asociada al evento $M$ (por ejemplo $\mathrm{xG}$, $\mathrm{xA}$, $\mathrm{xCards}$).
\end{itemize}

\section{El concepto xG y su generalización}
\label{sec:xg}

\subsection{Origen y motivación}

El concepto \emph{Expected Goals} (xG) fue introducido por Sam Green en 2012 en un blog de Opta, formalizando una intuición previa: no todos los tiros tienen la misma probabilidad de convertirse en gol. Un mano-a-mano con el arquero a tres metros de la portería y un disparo de treinta metros con cinco defensores en medio son eventos cualitativamente distintos, aunque ambos cuenten como ``un tiro'' en las estadísticas tradicionales.

La idea es asignar a cada tiro una probabilidad estimada de gol, derivada empíricamente de cuántas veces tiros con características similares se convirtieron en gol a lo largo de un corpus histórico. Esa probabilidad es la xG del tiro. Sumando xG sobre los tiros de un partido o temporada se obtiene la métrica agregada $\mathrm{xG}$ de un jugador o equipo.

\subsection{Definición formal}

Sea $x_i$ un tiro y $\mathbf{c}_i \in \mathbb{R}^k$ el vector de contexto observado en el momento del tiro. Sea $Y_i \in \{0,1\}$ el indicador de gol ($Y_i = 1$ si el tiro fue gol). Un modelo xG estima:

\[
  \mathrm{xG}(x_i) \;=\; P(Y_i = 1 \mid \mathbf{c}_i)
\]

\noindent
en la práctica vía regresión logística, \emph{gradient boosting} o redes neuronales. Para un conjunto de tiros $\{x_1,\dots,x_n\}$ se define la xG agregada como:

\[
  \mathrm{xG}_{\mathrm{agg}} \;=\; \sum_{i=1}^{n} \mathrm{xG}(x_i)
\]

\noindent
y la diferencia $G - \mathrm{xG}_{\mathrm{agg}}$, donde $G = \sum_i Y_i$ es el número de goles realmente convertidos, se interpreta como \emph{calidad de definición} (\emph{finishing skill}): valores positivos indican que el jugador convierte más de lo que la jugada promedio sugiere; valores negativos, lo contrario.

\subsection{Features típicas del vector de contexto}

Las variables que componen $\mathbf{c}_i$ varían entre modelos, pero suelen incluir:

\begin{table}[h]
\centering
\begin{tabular}{ll}
\toprule
\textbf{Feature} & \textbf{Interpretación} \\
\midrule
Distancia a la portería          & Variable más importante \\
Ángulo de tiro                   & Apertura geométrica del marco \\
Parte del cuerpo                 & Pie hábil vs no hábil vs cabeza \\
Tipo de jugada previa            & Pase, centro, rebote, contraataque \\
Bajo presión defensiva           & Defensores entre tirador y arco \\
Posición del arquero             & Sólo en modelos con \emph{freeze-frame} \\
Número de defensores             & Sólo en modelos con \emph{freeze-frame} \\
Velocidad del balón              & Sólo con datos de tracking \\
\bottomrule
\end{tabular}
\caption{Features comunes en modelos xG. Las dos últimas requieren datos de tracking de jugadores que sólo se exponen comercialmente.}
\end{table}

El modelo xG de StatsBomb publicado en su \emph{open data} usa el \emph{freeze-frame}: una fotografía del momento del tiro con la posición de todos los jugadores, lo que permite calcular cuántos defensores están entre el tirador y la portería, dónde está el arquero, etc. Otros modelos (FBref, Understat) usan sólo la geometría del tiro y son por lo tanto menos precisos.

\subsection{Interpretación práctica}

\begin{table}[h]
\centering
\begin{tabular}{lc}
\toprule
\textbf{Tipo de tiro} & \textbf{xG aproximado} \\
\midrule
Penal & 0.76 \\
\emph{Tap-in} a dos metros de portería vacía & 0.95 \\
Cabezazo desde el corazón del área & 0.20 -- 0.40 \\
Cabezazo desde primer palo, ángulo cerrado & 0.05 \\
Disparo desde fuera del área ($\sim$25 m) & 0.03 -- 0.05 \\
Tiro libre directo desde 25 m & 0.06 \\
Mano a mano con el arquero & 0.40 -- 0.60 \\
\bottomrule
\end{tabular}
\caption{Rangos típicos de xG según el tipo de tiro. Calibrados con datos abiertos de StatsBomb.}
\end{table}

\subsection{Por qué xG predice mejor que goles}

Los goles son eventos raros y de varianza muy alta partido a partido. Para un mismo nivel real de habilidad, un delantero puede meter tres goles en un partido y ninguno en los siguientes cuatro. xG es una métrica de varianza mucho menor: la cantidad de chances de calidad generadas por un jugador es relativamente estable y refleja con más fidelidad su contribución sistemática al juego.

Empíricamente, en nuestros datos abiertos de torneos internacionales (FIFA World Cup 2018, 2022; UEFA Euro 2020, 2024; Copa América 2024; \emph{etcétera}, total 908 partidos), la correlación de Pearson entre goles y xG es $r \approx 0.62$, mientras que entre goles y tiros (sin ajustar por calidad) es $r \approx 0.44$. Ese delta de $0.18$ es la mejora que aporta el ajuste por calidad de chance: xG es un predictor de goles futuros más fuerte que los goles propios pasados.

\subsection{Generalización a cualquier evento contable}

La construcción anterior se generaliza directamente. Para cualquier evento discreto $M$ que se pueda observar y etiquetar binariamente (\emph{ocurre / no ocurre}), si existe un vector de contexto $\mathbf{c}$ predictivo, se puede definir:

\[
  \mathrm{x}\textsc{M}(x_i) \;=\; P(M_i = 1 \mid \mathbf{c}_i)
\]

\noindent
y su agregado correspondiente. Esto da pie al catálogo descrito en la Sección~\ref{sec:catalogo}.

Los principios de identificación de un buen candidato para una métrica \emph{expected} son tres:

\begin{enumerate}[label=\textbf{P\arabic*.}]
  \item \textbf{Granularidad por evento}: cada ocurrencia debe ser identificable individualmente (no sólo un \emph{count} agregado).
  \item \textbf{Contexto observable}: deben existir features pre-evento que predigan razonablemente la ocurrencia.
  \item \textbf{Aplicabilidad a predicción}: el evento debe ser repetible suficientes veces para que el agregado tenga señal.
\end{enumerate}

\section{Catálogo de métricas \emph{expected}}
\label{sec:catalogo}

Esta sección enumera las métricas \emph{expected} más utilizadas en la literatura, agrupadas por dominio del juego. Para cada una se explicita su definición, los observables necesarios para calcularla y su interpretación práctica.

\subsection{Métricas ofensivas}

\subsubsection{xG --- Expected Goals}

\textbf{Definición.} Probabilidad de gol asignada a cada tiro, calibrada empíricamente.
\[
  \mathrm{xG}(x_i) = P(\text{gol} \mid \mathbf{c}_i)
\]
\textbf{Observables.} Evento de tiro con ubicación, parte del cuerpo, tipo de jugada previa.\\
\textbf{Interpretación.} Suma por jugador/equipo aproxima ``goles esperados''. La diferencia $G - \mathrm{xG}$ mide finalización.

\subsubsection{npxG --- non-penalty xG}

\textbf{Definición.} xG excluyendo los tiros desde penal.\\
\textbf{Motivación.} Los penales tienen xG fija ($\approx 0.76$) y dependen más del designado al penal que del juego abierto. npxG aísla la contribución al juego en movimiento.

\subsubsection{xA --- Expected Assists}

\textbf{Definición.} Para un pase $p_i$ que termina en tiro, $\mathrm{xA}(p_i) = \mathrm{xG}(\text{tiro resultante})$. Si no hay tiro, $\mathrm{xA}(p_i) = 0$.\\
\textbf{Observables.} Pasos completados que generan tiros (\emph{key passes}), más la xG de esos tiros.\\
\textbf{Interpretación.} Cuantifica creación de chances. Mejor que asistencias crudas porque no penaliza al pasador cuando el rematador falla.

\subsubsection{xAG --- Expected Assisted Goals}

\textbf{Definición.} Variante de xA que sólo cuenta el pase inmediatamente previo a un tiro, no toda la cadena.

\subsubsection{xT --- Expected Threat}

\textbf{Definición.} Mapa de valor sobre la cancha que asigna a cada celda $(x,y)$ una probabilidad esperada de que la posesión termine en gol dentro de los próximos $N$ eventos. La diferencia $\Delta \mathrm{xT}$ entre origen y destino de un pase o carrera valora la jugada.\\
\textbf{Observables.} Cadena de eventos por posesión con coordenadas.\\
\textbf{Interpretación.} Mide ``progresión peligrosa''. Un centrocampista que avanza balón 30 metros hacia el área genera $\Delta \mathrm{xT}$ alto sin necesidad de generar tiro.

\subsubsection{xPass --- Expected Pass Completion}

\textbf{Definición.} Probabilidad de que un pase se complete, dada distancia, ángulo y presión.\\
\textbf{Uso.} Ajustar el porcentaje de completitud por dificultad. Un centrocampista que completa $75\%$ de pases difíciles vale más que uno que completa $95\%$ de pases laterales triviales.

\subsection{Métricas defensivas y de presión}

\subsubsection{xGA --- Expected Goals Against}

\textbf{Definición.} xG agregada de los tiros recibidos por un equipo. Métrica defensiva paralela a xG.

\subsubsection{PPDA --- Passes per Defensive Action}

\textbf{Definición.} Pases permitidos al rival en el tercio defensivo propio antes de cometer una acción defensiva (entrada, intercepción, falta).\\
\textbf{Interpretación.} Cuantifica intensidad de presión: valores bajos indican presión alta y agresiva.

\subsection{Métricas disciplinarias}

\subsubsection{xFouls --- Expected Fouls Committed}

\textbf{Definición.} Probabilidad de que un duelo / acción defensiva termine en falta, calibrada por contexto (zona, estilo del jugador, árbitro).\\
\textbf{Uso.} Identificar jugadores que se exponen sistemáticamente a faltas en zonas peligrosas.

\subsubsection{xFD --- Expected Fouls Drawn}

\textbf{Definición.} Probabilidad de que un drible / recepción de un jugador termine en falta a favor.\\
\textbf{Interpretación.} Habilidad para forzar faltas. Mide tanto skill (\emph{regateadores} provocan más faltas) como riesgo.

\subsubsection{xCards --- Expected Cards}

\textbf{Definición.} Probabilidad de que un jugador sea amonestado en un partido, dada su tasa histórica, el estilo del árbitro y el contexto del encuentro (rivalidad, importancia, marcador).\\
\textbf{Observables.} Eventos \texttt{Foul Committed} con campo \texttt{card}, eventos \texttt{Bad Behaviour}, asignación de árbitro.\\
\textbf{Uso típico.} Mercados de apuestas; análisis disciplinario; ajuste de modelos de minutos esperados.

\subsection{Métricas de pelota detenida (\emph{set pieces})}

\subsubsection{xCorners --- Expected Corners}

\textbf{Definición.} Corners esperados por equipo en un partido. Función de \emph{shots blocked}, presión ofensiva y estilo defensivo del rival.\\
\textbf{Observables.} Eventos con \texttt{pass\_type = Corner} (o equivalentes); cantidad de tiros desviados / despejados.

\subsubsection{xFreeKicks --- Expected Free Kicks}

\textbf{Definición.} Tiros libres esperados (sumando todas las posiciones, o restringido al último tercio).\\
\textbf{Uso.} Modelar oportunidades de \emph{set piece} ofensiva; identificar jugadores que las generan.

\subsubsection{xThrows --- Expected Throw-Ins}

\textbf{Definición.} Saques de banda esperados por equipo en un partido. Mecánico pero útil como predictor de duración efectiva.\\
\textbf{Limitación.} Métrica de bajo \emph{insight} por sí misma; su valor está en modelos de mercado especializado.

\subsubsection{xGoalKicks --- Expected Goal Kicks}

\textbf{Definición.} Saques de meta esperados; refleja tiros desviados y presión ofensiva sostenida.

\subsection{Métricas compositas}

\subsubsection{VAEP --- Valuing Actions by Estimated Probabilities}

\textbf{Definición.} Framework de Decroos y otros (2019) que asigna a cada acción un valor en goles esperados: $V(a_i) = \Delta P(\text{gol propio}) - \Delta P(\text{gol rival})$, considerando los efectos a corto plazo de la acción.\\
\textbf{Diferencia con xT.} VAEP también valora acciones defensivas (despejes, intercepciones), no sólo progresión ofensiva.

\subsubsection{G+ / Goals Added}

\textbf{Definición.} Métrica de la MLS análoga a VAEP. Diferencias de implementación, mismo principio.

\subsection{Tabla resumen}

\begin{longtable}{lp{5.5cm}p{4cm}p{2.5cm}}
\toprule
\textbf{Métrica} & \textbf{Mide} & \textbf{Requiere eventos de\ldots} & \textbf{Granularidad mínima} \\
\midrule
\endhead
xG       & calidad de chance de un tiro             & tiros con coordenadas      & por tiro \\
npxG     & xG sin penales                            & tiros + tipo               & por tiro \\
xA       & creación de chance                        & pases + tiros consecutivos & por pase \\
xT       & valor de mover el balón                   & cadena de eventos          & por evento \\
xPass    & dificultad de un pase                     & pases con coordenadas      & por pase \\
xGA      & calidad de chance concedida               & tiros del rival            & por tiro \\
PPDA     & intensidad de presión                     & pases del rival + defensa  & por equipo-partido \\
xFouls   & probabilidad de cometer falta             & faltas con contexto        & por evento \\
xFD      & probabilidad de recibir falta             & faltas con contexto        & por evento \\
xCards   & probabilidad de amonestación              & tarjetas + árbitro         & por jugador-partido \\
xCorners & corners esperados                         & corners + flujo ofensivo   & por equipo-partido \\
xFK      & tiros libres esperados                    & faltas con ubicación       & por equipo-partido \\
xThrows  & saques de banda esperados                 & restarts                   & por equipo-partido \\
xGK      & saques de meta esperados                  & restarts del rival         & por equipo-partido \\
VAEP     & valor en goles esperados de una acción   & todos los eventos          & por evento \\
\bottomrule
\caption{Catálogo resumido. Granularidad mínima indica el nivel al que típicamente se reporta la métrica.}
\end{longtable}

\section{Fuentes de datos y cobertura}
\label{sec:fuentes}

Las métricas \emph{expected} sólo son construibles si la fuente expone los \emph{observables} requeridos. Esta sección caracteriza las fuentes públicas más usadas y cruza su disponibilidad contra el catálogo de la Sección~\ref{sec:catalogo}.

\subsection{Fuentes públicas relevantes}

\subsubsection{StatsBomb Open Data}

\textbf{Forma de acceso}: archivos JSON publicados directamente en GitHub (\url{https://github.com/statsbomb/open-data}). No hay scraping de páginas web ni rate limit relevante.\\
\textbf{Granularidad}: \emph{event-level}. Cada partido se expone como una lista de $\sim 3000$ eventos, cada uno con tipo (Pass, Shot, Foul, Carry, etc.), coordenadas $(x,y)$, tiempo y atributos específicos del tipo. Cada tiro incluye \emph{freeze-frame} con la posición de todos los jugadores en el momento.\\
\textbf{Cobertura}: limitada a competiciones específicas que StatsBomb decide liberar. En orden de importancia ofensiva para el proyecto: FIFA World Cup 2018, 2022; UEFA Euro 2020, 2024; UEFA Women's Euro 2022, 2025; Copa América 2024; Champions League 2003/04--2018/19; Premier League 2003/04 y 2015/16; La Liga 2004/05--2020/21 (algunas restringidas a partidos de Messi); Major League Soccer 2023; Indian Super League 2021/22; FA WSL varias temporadas.\\
\textbf{Licencia}: open con atribución requerida.

\subsubsection{FBref / Sports Reference}

\textbf{Forma de acceso}: scraping HTML de \url{https://fbref.com}. Sitio detrás de Cloudflare con \emph{JS challenge}, requiere navegador real (Chromium). Rate limit recomendado: 1 request cada 3 segundos.\\
\textbf{Granularidad}: \emph{aggregate} y \emph{match summary}. No expone eventos individuales. A nivel temporada, da agregados por jugador (\texttt{standard}, \texttt{shooting}, \texttt{misc}, \texttt{playing\_time}, \texttt{keeper}). A nivel partido, da el \emph{summary} con counts agregados de tiros, tackles, intercepciones, faltas y tarjetas.\\
\textbf{Cobertura}: amplísima --- Big~5 europeas, MLS, ligas sudamericanas, copas internacionales, históricos hasta los 90.\\
\textbf{Fuente original}: datos licenciados de Opta/StatsPerform.

\subsubsection{Understat}

\textbf{Forma de acceso}: scraping HTML simple de \url{https://understat.com}; sin Cloudflare.\\
\textbf{Granularidad}: \emph{shot-level}. Cada tiro con coordenadas y xG calculado por modelo propio. Agregables a nivel partido y temporada.\\
\textbf{Cobertura}: Big~5 europeas + RFPL desde 2014. Cobertura continua, sin gaps.\\
\textbf{Limitación}: foco en tiros y xG; no expone pases, dribles, faltas con detalle.

\subsubsection{ClubElo}

\textbf{Forma de acceso}: API simple en \url{http://clubelo.com}; CSVs y JSON.\\
\textbf{Granularidad}: rating diario por club.\\
\textbf{Uso}: feature de \emph{strength of opponent} en modelos jugador-partido. Disponible para todos los clubes europeos relevantes con histórico de décadas.

\subsubsection{WhoScored / Sofascore}

\textbf{Forma de acceso}: scraping pesado en JavaScript; soccerdata lo soporta vía Chromium.\\
\textbf{Granularidad}: \emph{match-level} con stats granulares por jugador (incluye ratings propios).\\
\textbf{Cobertura}: amplia.\\
\textbf{Limitación}: scraping lento y vulnerable a cambios de la web.

\subsubsection{Transfermarkt}

\textbf{Forma de acceso}: scraping HTML simple, con paginación pero sin Cloudflare agresivo.\\
\textbf{Contenido}: \emph{no} expone estadísticas de partido; sí expone metadata de jugadores (fecha de nacimiento, altura, pie hábil, valor de mercado, historial de clubes, lesiones).\\
\textbf{Uso}: enriquecer el catálogo \texttt{players.csv}.

\subsubsection{APIs comerciales}

\textbf{API-Football} (RapidAPI): cobertura global, free tier 100 req/día, pago desde \$20/mes. Incluye eventos por partido (no event-level rico, sí counts).\\
\textbf{Sportmonks}, \textbf{Sportradar}, \textbf{Opta}: comerciales, fuera de scope para uso personal.

\subsection{Matriz métrica $\times$ fuente}

Las celdas indican factibilidad y comentarios:

\begin{itemize}[noitemsep]
  \item \textbf{Sí}: la fuente expone los observables directamente.
  \item \textbf{Derivable}: posible reconstruir desde los datos crudos con trabajo.
  \item \textbf{No}: la fuente no contiene los observables suficientes.
\end{itemize}

\begin{table}[h]
\centering
\small
\begin{tabular}{lcccc}
\toprule
\textbf{Métrica} & \textbf{StatsBomb} & \textbf{FBref} & \textbf{Understat} & \textbf{WhoScored} \\
\midrule
xG por tiro                 & Sí (con freeze-frame) & No directo & Sí (propio)   & Sí \\
xG agregada jugador-temporada & Derivable           & Sólo Big 5\,$^{*}$ & Sí       & Sí \\
xG por jugador-partido       & Derivable             & Derivable & Derivable      & Sí \\
xA por jugador-partido       & Derivable             & Sí parcial & Limitado      & Sí \\
xT                           & Derivable             & No        & No             & No \\
xCards / xFouls              & Derivable             & Counts no expected & No   & Counts \\
xCorners                     & Derivable             & Counts (team-match) & No  & Sí \\
xThrows / xGoalKicks         & Derivable             & No        & No             & No \\
PPDA                         & Derivable             & No        & No             & Parcial \\
VAEP                         & Derivable             & No        & No             & No \\
\bottomrule
\end{tabular}
\caption{Matriz métrica $\times$ fuente. ``Derivable'' implica que el script de agregación es responsabilidad nuestra, no de la fuente. \\$^{*}$ FBref Big 5 \emph{Combined} omite el bloque \emph{Expected} en el HTML que scrapea soccerdata (limitación upstream).}
\end{table}

\subsection{Estrategia de cobertura del proyecto}

Dado el catálogo y las fuentes, la estrategia es:

\begin{enumerate}[label=\textbf{S\arabic*.}]
  \item \textbf{StatsBomb} es la fuente primaria para todo lo \emph{expected} en las 9 competiciones cubiertas (xG, xA, xT, xFouls, xCards, xCorners, etc.). 3.1 millones de eventos en disco, casi 26 mil filas de jugador-partido agregado.
  \item \textbf{FBref} es la fuente primaria para counts \emph{aggregate} en Big 5 a nivel temporada (5 temporadas, 17.5 mil filas jugador-temporada).
  \item \textbf{FBref match summary} es la fuente para counts \emph{per-match} en Big 5 actual (lo que se va construyendo en la Fase A del roadmap).
  \item \textbf{Understat} llena el gap de xG para Big 5 a nivel temporada (FBref no lo expone vía soccerdata).
  \item \textbf{ClubElo} aporta la feature de \emph{strength of opponent} requerida para modelos de jugador-partido bien especificados.
  \item \textbf{Transfermarkt} se considera para enriquecer \texttt{players.csv} con metadata biográfica.
\end{enumerate}

\subsection{Resumen ejecutivo: ¿de dónde sale cada métrica para este proyecto?}

\begin{table}[h]
\centering
\small
\begin{tabular}{lll}
\toprule
\textbf{Métrica objetivo} & \textbf{Fuente primaria} & \textbf{Cobertura efectiva} \\
\midrule
xG, npxG por tiro                 & StatsBomb         & 9 competiciones, 908 partidos \\
xG por jugador-temporada Big 5     & Understat (planificado) & Big 5, 11 temporadas \\
xA, xT, xPass, VAEP                & StatsBomb         & 9 competiciones \\
xFouls / xFouls Drawn              & StatsBomb         & 9 competiciones \\
xCards                             & StatsBomb (+ árbitro de FBref) & Limitado \\
xCorners / xThrows / xGoalKicks    & StatsBomb         & 9 competiciones \\
Counts tiros, tackles, etc.\ Big 5 & FBref match summary & Big 5 partidos scrapeados \\
Counts jugador-temporada Big 5     & FBref aggregates  & Big 5, 5 temporadas \\
Strength of opponent               & ClubElo (planificado) & Todos los equipos europeos \\
Metadata biográfica de jugadores   & Transfermarkt (futuro) & Universal \\
\bottomrule
\end{tabular}
\caption{Asignación final fuente $\to$ métrica adoptada en \textsc{Futbolstats}.}
\end{table}

\section{Implementación}
\label{sec:impl}

Esta sección documenta cómo está cableado el repositorio para producir las tablas de las que parten los modelos.

\subsection{Capas del pipeline}

El pipeline tiene cuatro capas con responsabilidades claras:

\begin{enumerate}[label=\textbf{L\arabic*.}]
  \item \textbf{Ingestion} (\texttt{06\_ingestion/}): scrapers y \emph{loaders} que persisten datos crudos en \texttt{01\_data\_raw/} con caché HTML/JSON.
  \item \textbf{Entities} (\texttt{03\_entities/}): catálogos maestros \texttt{players.csv}, \texttt{teams.csv}, \texttt{countries.csv}, \texttt{competitions.csv} con IDs canónicos prefijados por fuente (\texttt{fbref\_<id>}, \texttt{sb\_<id>}).
  \item \textbf{Features} (\texttt{07\_features/}): scripts que toman crudos y producen tablas planas listas para ML en \texttt{02\_data\_processed/}.
  \item \textbf{Models} (\texttt{08\_models/}) y \textbf{Correlations} (\texttt{09\_correlations/}): consumen las tablas procesadas.
\end{enumerate}

\subsection{Scripts por fuente}

\subsubsection{FBref}

\begin{itemize}[noitemsep]
  \item \texttt{fbref\_scraper.py}: stats agregadas por jugador-temporada, 5 \emph{stat types} disponibles vía soccerdata (\texttt{standard}, \texttt{shooting}, \texttt{playing\_time}, \texttt{misc}, \texttt{keeper}).
  \item \texttt{fbref\_match\_scraper.py}: stats por jugador-partido (\texttt{summary} y \texttt{keepers}). Una request HTTP por partido (caché HTML reutilizable).
\end{itemize}

\subsubsection{StatsBomb}

\begin{itemize}[noitemsep]
  \item \texttt{statsbomb\_loader.py}: descarga matches, events, lineups por competición. CLI con \texttt{-{}-list}, \texttt{-{}-competition}, \texttt{-{}-season}.
  \item \texttt{statsbomb\_aggregate.py}: roll-up por jugador-partido de las stats principales (goles, asistencias, xG, pases, tackles, intercepciones, faltas, dribles).
  \item \texttt{statsbomb\_entities.py}: \emph{append-only} de jugadores y equipos a los catálogos.
\end{itemize}

\subsubsection{Understat}

\begin{itemize}[noitemsep]
  \item \texttt{understat\_scraper.py}: stats agregadas por jugador-temporada y por jugador-partido, con xG calculado por modelo propio de Understat.
\end{itemize}

\subsubsection{ClubElo}

\begin{itemize}[noitemsep]
  \item \texttt{clubelo\_scraper.py}: ratings diarios por club; usado como feature de fuerza del rival.
\end{itemize}

\subsection{Scripts de \emph{features}}

\begin{itemize}[noitemsep]
  \item \texttt{build\_player\_season.py}: une las 5 \emph{stat types} de FBref por (player, team, league, season). Output: \texttt{fbref\_Big5\_player\_season.parquet} ($\sim 17\,500$ filas, 85 columnas).
  \item \texttt{build\_player\_match.py}: enriquece el \texttt{summary} de FBref con el \texttt{schedule} (fecha, local/visitante, oponente). Output: \texttt{fbref\_<comp>\_player\_match.parquet}.
  \item \texttt{build\_discipline\_setpieces.py}: agrega cards, fouls, corners, throw-ins, goal kicks y free kicks desde los eventos StatsBomb. Outputs: \texttt{statsbomb\_player\_match\_discipline.parquet}, \texttt{statsbomb\_team\_match\_setpieces.parquet}.
\end{itemize}

\subsection{Schemas de los outputs principales}

\subsubsection{\texttt{fbref\_Big5\_player\_season.parquet}}

Una fila por jugador y temporada. Identidad: \texttt{(player, team, league, season)}.\\
Columnas con sufijo \texttt{\_\_<stat>}: \texttt{Performance\_Gls\_\_standard}, \texttt{Standard\_Sh\_\_shooting}, \texttt{Performance\_CrdY\_\_misc}, etc.

\subsubsection{\texttt{statsbomb\_player\_match.parquet}}

Una fila por jugador y partido. 28 columnas: \texttt{player\_id, match\_id, team\_id, opponent\_id, competition\_id, season, date, is\_home, minutes\_played, goals, assists, shots, shots\_on\_target, xg, passes\_attempted, passes\_completed, key\_passes, dribbles\_attempted, dribbles\_completed, tackles, interceptions, fouls\_committed, fouls\_drawn, \ldots}

\subsubsection{\texttt{statsbomb\_player\_match\_discipline.parquet}}

Una fila por jugador y partido. Cards y fouls desagregados:
\begin{itemize}[noitemsep]
  \item \texttt{yellow\_cards}, \texttt{second\_yellows}, \texttt{red\_cards}
  \item \texttt{fouls\_committed}, \texttt{fouls\_drawn}
  \item \texttt{fouls\_committed\_def\_third}, \texttt{fouls\_committed\_mid\_third}, \texttt{fouls\_committed\_att\_third}
\end{itemize}

\subsubsection{\texttt{statsbomb\_team\_match\_setpieces.parquet}}

Una fila por equipo y partido. Counts de \emph{set pieces}:
\begin{itemize}[noitemsep]
  \item \texttt{corners\_for}, \texttt{free\_kicks\_for}, \texttt{throw\_ins\_for}, \texttt{goal\_kicks\_for}
  \item Mismo set con \texttt{\_against} para los recibidos.
\end{itemize}

\subsection{Flujo de procesamiento}

\begin{verbatim}
                +-----------------+
                |  Fuentes web    |
                | (FBref, SB, US, |
                |  ClubElo, etc.) |
                +--------+--------+
                         |
                         v
                +-----------------+
                |  06_ingestion/  |
                | (scrapers,      |
                |  loaders)       |
                +--------+--------+
                         |
                         v
                +-----------------+
                |  01_data_raw/   |
                | (parquet + HTML |
                |  cache)         |
                +--------+--------+
                         |
              +----------+----------+
              |                     |
              v                     v
     +-----------------+   +-----------------+
     |  03_entities/   |   |  07_features/   |
     | (players.csv,   |   | (build_*.py)    |
     |  teams.csv...)  |   +--------+--------+
     +-----------------+            |
                                    v
                          +-------------------+
                          | 02_data_processed |
                          | (parquet, ML-     |
                          |  ready tables)    |
                          +---------+---------+
                                    |
                                    v
                          +-------------------+
                          | 08_models/        |
                          | 09_correlations/  |
                          +-------------------+
\end{verbatim}

\subsection{Cómo extender con una nueva fuente}

\begin{enumerate}[label=\textbf{E\arabic*.}]
  \item Agregar \texttt{06\_ingestion/<source>\_loader.py} siguiendo el patrón existente: salida en \texttt{01\_data\_raw/<entity>/<source>\_<scope>.parquet}, caché en \texttt{01\_data\_raw/\_cache/<source>/}.
  \item Usar \texttt{lib.io.write\_parquet} (compresión zstd nivel 9 enforzada).
  \item Prefijar todos los IDs nuevos con \texttt{<source>\_} (por ejemplo \texttt{us\_} para Understat).
  \item Si la fuente requiere \emph{append} a los catálogos maestros, crear \texttt{<source>\_entities.py} análogo al de StatsBomb.
  \item Documentar quirks y rate limits en \texttt{.claude/skills/<source>-ingest/SKILL.md}.
\end{enumerate}

\section{Pipeline de \emph{value betting}}
\label{sec:betting}

Esta sección documenta la capa que conecta las predicciones de estadísticas
individuales con las cuotas de las casas de apuestas, para detectar apuestas de
valor esperado positivo (\emph{+EV}) sobre \emph{player props}.

\subsection{De la predicción a la decisión}

El flujo convierte una predicción del modelo en una decisión de apuesta en cuatro
pasos:

\[
\underbrace{\lambda}_{\text{predicción}}
\;\longrightarrow\;
\underbrace{P(Y \gtrless \ell)}_{\text{prob.\ over/under}}
\;\longrightarrow\;
\underbrace{\mathrm{EV} = p\cdot d - 1}_{\text{valor esperado}}
\;\longrightarrow\;
\underbrace{f^\star = \tfrac{p\,b - q}{b}}_{\text{Kelly}}
\]

donde $\lambda$ es la media predicha del conteo (tiros, tarjetas, etc.), $\ell$ la
línea de la casa, $d$ la cuota decimal, $b = d-1$ las cuotas netas, $p$ la
probabilidad del modelo y $q = 1-p$.

\subsection{Modelos de conteo para over/under}

Una apuesta \emph{over/under} sobre una estadística de conteo requiere
$P(Y > \ell)$ y $P(Y < \ell)$. Modelamos $Y$ como variable de conteo:

\begin{itemize}[noitemsep]
  \item \textbf{Poisson}$(\lambda)$: supone $\mathrm{Var}(Y) = \mathbb{E}(Y) = \lambda$.
  \item \textbf{Binomial Negativa} con media $\mu$ y dispersión $\alpha$, tal que
  $\mathrm{Var}(Y) = \mu + \alpha\mu^2$. Se reduce a Poisson cuando $\alpha \to 0$.
  Conversión a la parametrización $(n,p)$ de \texttt{scipy}: $n = 1/\alpha$,
  $p = 1/(1+\alpha\mu)$.
\end{itemize}

Las estadísticas reales de fútbol (tiros, tackles, tarjetas) están casi siempre
\emph{sobre-dispersas} ($\mathrm{Var} > \text{media}$), por lo que la Binomial
Negativa es el modelo honesto. El parámetro $\alpha$ se estima por momentos:
$\hat{\alpha} = (\mathrm{Var} - \text{media}) / \text{media}^2$. Para líneas
semienteras (p.\,ej.\ $\ell = 2.5$) no hay empate; para líneas enteras
($\ell = 2$) el resultado exacto produce un \emph{push} (devolución del stake).

\subsection{Eliminación del \emph{vig}}

La probabilidad implícita de una cuota, $1/d$, incluye el margen de la casa: la
suma sobre los resultados de un mercado (\emph{booksum}) excede 1 por el
\emph{overround}. Para comparar contra el modelo se devig\'ua. El método
proporcional divide cada implícita por el \emph{booksum}; el método de potencia
resuelve el exponente $k$ tal que $\sum_i p_i^{\,k} = 1$, reduciendo el sesgo
favorito--\emph{longshot}. Ambos en \texttt{lib.betting.remove\_vig}.

\subsection{Dimensionamiento por Kelly}

El criterio de Kelly maximiza el crecimiento logarítmico del bankroll apostando
una fracción $f^\star = (p\,b - q)/b$. En la práctica se usa \emph{Kelly
fraccional} (un cuarto, $f = 0.25\,f^\star$) porque Kelly completo asume $p$
conocido con certeza; con $p$ estimado, el Kelly completo produce \emph{drawdowns}
brutales. La incertidumbre sobre la calibración del modelo justifica el
descuento.

\subsection{Reorganización de la base de datos}

La tabla de entrenamiento (\texttt{02\_data\_processed/events\_clean/}) es plana,
una fila por jugador-partido. El análisis de apuestas necesita otra forma,
centrada en \emph{partidos $\to$ jugadores $\to$ mercados $\to$ cuotas $\to$
valor}. El script \texttt{07\_features/build\_matches\_view.py} construye esa vista
en \texttt{02\_data\_processed/matches\_view/}:

\begin{itemize}[noitemsep]
  \item \texttt{fixtures.parquet}: un partido por fila (fecha, liga, local/visitante, estado).
  \item \texttt{match\_players.parquet}: por (partido, jugador), predicción
  $\lambda$ \emph{leak-safe} y resultado realizado de cada mercado.
  \item \texttt{match\_odds.parquet}: cuotas en formato largo (lo llena el scraper).
  \item \texttt{match\_value.parquet}: el \emph{join} con \texttt{p\_model},
  \texttt{implied\_prob}, \texttt{edge}, \texttt{ev}, \texttt{kelly} y el resultado
  para backtesting.
\end{itemize}

La predicción $\lambda$ por defecto es el baseline \emph{leak-safe}
(media móvil de 5 partidos con \texttt{shift(1)}); se sustituye por predicciones de
XGBoost sobre-escribiendo las columnas \texttt{lam\_<market>}.

\subsection{Fuente de cuotas: Pinnacle}

\texttt{06\_ingestion/pinnacle\_scraper.py} consume la API guest JSON de Pinnacle
(\texttt{guest.api.arcadia.pinnacle.com}), sin Selenium ni Cloudflare. Se eligió
Pinnacle por su bajo margen, límites altos y por no vetar ganadores, lo que la
hace el mejor \emph{benchmark} de probabilidad de mercado. Limitación importante:
para fútbol Pinnacle sólo cotiza \textbf{props de goles y tarjetas} (anytime/first
goalscorer, \emph{to score}, \emph{to be booked}); no ofrece over/under de
tiros/tackles/pases. El edge sobre esos mercados requiere un book que los cotice
(p.\,ej.\ Bet365).

\subsection{Validación y \emph{backtest}}

\texttt{lib.backtest} y \texttt{08\_models/backtest\_betting.py} simulan la
estrategia: PnL por unidad $= (d-1)$ si gana, $-1$ si pierde, $0$ en \emph{push}.
Se evalúa el ROI según el umbral de EV, con stake plano y Kelly fraccional. La
\textbf{calibración} —¿cuando el modelo dice $p=0.6$ acierta el 60\,\%?— es el
guardarraíl: si no calibra, el EV miente. El \emph{backtest} de validación corre
sobre un mercado sintético cotizado por posición; su ROI positivo verifica que la
\emph{lógica} funciona, pero \textbf{no es una afirmación de edge real}: contra un
book sharp el EV esperado es $\approx -\text{vig}$.

\section{Limitaciones y trabajo futuro}
\label{sec:limitaciones}

Esta sección discute los principales puntos débiles del pipeline actual y caminos para extenderlo.

\subsection{Cobertura desigual entre fuentes}

La cobertura de StatsBomb \emph{Open Data} es excelente en profundidad (eventos individuales con \emph{freeze-frame}) pero limitada en amplitud: cubre torneos internacionales y un puñado de temporadas-liga seleccionadas. No hay cobertura continua de Big~5 europeas con event data abiertos.

FBref, en contraste, tiene cobertura continua de Big~5 y otras ligas, pero solo a nivel agregado o de \emph{match summary} (counts), sin eventos individuales. La consecuencia práctica es asimétrica:

\begin{itemize}[noitemsep]
  \item Modelos de jugador-partido para Copa~América, World Cup, Champions League (datos StatsBomb): se pueden construir con todas las métricas \emph{expected} del catálogo.
  \item Modelos de jugador-partido para Premier League 2024-25: se construyen con counts de FBref \texttt{summary} y agregados StatsBomb donde haya overlap; xCorners y xThrows requerirían \emph{otra} fuente (WhoScored, API-Football pago).
\end{itemize}

\subsection{xG en Big~5 vía FBref}

soccerdata, al scrapear la página \emph{Big 5 European Leagues Combined / shooting}, no recupera el \emph{supergroup} ``Expected'' (xG, npxG, etc.). La inspección del HTML cacheado lo confirma: el HTML servido a soccerdata carece del bloque. Si bien las páginas por liga individual sí lo tienen, el wrapper actual no las consulta.

\textbf{Mitigación elegida.} Agregar Understat como fuente complementaria. Understat publica su propio modelo xG para Big~5 + RFPL desde 2014; no es idéntico al de StatsBomb pero está bien calibrado y es ampliamente usado en la literatura.

\textbf{Caveat.} Los xG de Understat y StatsBomb \emph{no son intercambiables} para un mismo tiro: usan features distintas (Understat no tiene \emph{freeze-frame}). Para análisis longitudinal usar siempre la misma fuente.

\subsection{Sesgos de los modelos xG y derivados}

\subsubsection{Calidad del rematador}

xG \emph{no} incluye al rematador como feature. Esto es deliberado: si el modelo se enterara de quién patea, la diferencia $G - \mathrm{xG}$ dejaría de medir finalización individual. Pero implica que xG no es una estimación insesgada de la probabilidad real de gol cuando el rematador es atípico (Messi, Mbappé). Esos jugadores \emph{deben} estar arriba de su xG sistemáticamente.

\subsubsection{Sesgo del rival y del entorno}

Los modelos xG abiertos no usan al equipo defensor ni al arquero como features. Esto inflará la xG contra arqueros muy buenos y la deflará contra defensas que conceden sólo tiros lejanos. Para modelos individuales este sesgo es chico; para modelos de \emph{partido} es relevante y motiva el uso de ClubElo o equivalentes como feature de contexto.

\subsection{xCards y la variable árbitro}

Modelar \emph{expected cards} sin la identidad del árbitro deja fuera buena parte de la varianza: los árbitros difieren enormemente en su umbral de amonestación. StatsBomb \emph{Open Data} no expone el árbitro en sus archivos públicos. FBref sí lo tiene en su tabla de \texttt{schedule}.

\textbf{Mitigación futura.} Joinear \texttt{statsbomb\_player\_match\_discipline.parquet} con \texttt{fbref\_<comp>\_schedule\_<season>.parquet} cuando exista overlap (limitado, dada la asimetría de cobertura).

\subsection{Sesgo de selección en eventos defensivos}

Eventos defensivos como \texttt{Tackle}, \texttt{Interception} o \texttt{Block} están sujetos a sesgo de selección: un equipo que tiene posesión todo el partido no tiene oportunidad de hacer interceptaciones. Las métricas \emph{expected} sobre estos eventos deberían ajustar por exposición (minutos defendiendo, pases del rival permitidos en zona), o reportarse como \emph{tasa por exposición}, no como count crudo.

\subsection{Identificación de jugadores entre fuentes}

soccerdata no expone los \texttt{player\_id} hex de FBref. Esto fuerza al proyecto a sintetizar IDs por \texttt{(player\_name, team)}, lo que rompe cuando un jugador cambia de equipo a mitad de temporada o cuando dos jugadores comparten nombre. La solución correcta es scrapear directamente la página de jugador de FBref (que sí expone el hex ID) y mantener una tabla de \emph{mappings} entre fuentes.

\subsection{Splits temporales y \emph{data leakage}}

Todos los modelos del proyecto deben usar splits temporales (entrenar en temporada $T-1$, validar en $T$). Aleatorizar filas es \emph{data leakage} grueso: un jugador presente en train y test simultáneamente vuelve la métrica de error optimista en órdenes de magnitud.

Adicionalmente, las \emph{features} de \emph{rolling form} (\emph{e.g.}, ``goles últimos 5 partidos'') deben calcularse con un \texttt{shift(1)} para no incluir el partido actual. El script \texttt{shots\_baseline.py} ya implementa esta convención y debe usarse como plantilla.

\subsection{Líneas de trabajo futuro}

\begin{enumerate}[label=\textbf{F\arabic*.}]
  \item \textbf{xCards condicionado por árbitro}: scrapear referees desde FBref \texttt{schedule}, joinear con eventos StatsBomb, entrenar un modelo con árbitro como efecto fijo / aleatorio.
  \item \textbf{Modelos jerárquicos bayesianos}: efectos parcialmente \emph{pooled} por jugador, equipo y temporada. Implementables con PyMC; mejoran cuando hay poca data por entidad.
  \item \textbf{xT con datos públicos}: construir un modelo de \emph{expected threat} con la grilla de cancha de Karun Singh sobre eventos StatsBomb.
  \item \textbf{Enriquecimiento biográfico}: scraper de Transfermarkt para edad, altura, pie hábil, valor de mercado de cada jugador.
  \item \textbf{Splits de prueba más severos}: \emph{cold start} (jugador sin historia en train), \emph{cross-league} (entrenar en Big~5 y validar en torneos internacionales).
  \item \textbf{Embeddings de jugadores y equipos}: representaciones aprendidas que se compartan entre modelos de stats distintas.
  \item \textbf{Calibración}: comparar la distribución de predicciones contra la realidad. Una métrica $\hat{Y}$ está bien calibrada si, al promediar sobre todas las predicciones con $\hat{Y} \approx p$, el promedio empírico de $Y$ también es $\approx p$.
\end{enumerate}


\end{document}
